\chapter{Bell's Palsy}
\index{Bell's Palsy}

\section*{Definition and Overview}
Bell's palsy is a sudden weakness or paralysis of the muscles on one side of the face, caused by temporary dysfunction of the facial nerve (the seventh cranial nerve). The facial nerve supplies the muscles of facial expression, including those that raise the eyebrow, close the eye, and smile, as well as parts of the ear and tongue. The condition typically develops over hours to a few days, with maximal weakness usually reached within 48 to 72 hours.

Bell's palsy is a diagnosis of exclusion: it is the label applied to an acute, unilateral, lower-motor-neurone facial palsy for which no specific cause is found. The great majority of people recover fully or nearly fully, often within weeks to a few months, and the condition rarely leaves permanent disability. It is important to distinguish Bell's palsy from stroke and other more serious causes of facial weakness, which is why prompt medical assessment is always recommended.

\section*{Epidemiology}
Bell's palsy is the most common cause of one-sided facial paralysis, affecting roughly 15 to 30 people per 100,000 each year.

\begin{itemize}
    \item \textbf{Age distribution:} most cases occur between the ages of 15 and 60, with a peak in the 40s
    \item \textbf{Sex:} men and women are affected about equally
    \item \textbf{Pregnancy:} it is more common in pregnant women, particularly in the third trimester and in the first week after delivery
    \item \textbf{Risk factors:} diabetes, hypertension, obesity, and a recent upper respiratory tract infection are associated with a higher risk
    \item \textbf{Recurrence:} around 8 to 12 per cent of people experience a recurrence, sometimes on the opposite side, often years later
\end{itemize}

\section*{Aetiology and Pathophysiology}
The exact cause of Bell's palsy remains incompletely understood, but the leading theory is that viral infection or reactivation of a latent virus triggers inflammation and swelling of the facial nerve. Herpes simplex virus type 1 (the cold sore virus) is the most frequently implicated organism, although varicella zoster virus, cytomegalovirus, and Epstein--Barr virus have also been proposed.

The facial nerve travels through a long, narrow bony canal (the facial canal) within the temporal bone of the skull. When the nerve becomes inflamed and swollen inside this tight channel, it is compressed, its microvascular blood supply is compromised, and the conduction of nerve impulses is interrupted. The result is weakness or paralysis of the muscles supplied by the nerve. Because the nerve fibres are usually not permanently destroyed, the condition is typically reversible once the inflammation settles and the swelling subsides. Recovery depends on the degree of nerve injury, which ranges from mild compression (neurapraxia) to more severe damage requiring regeneration of nerve fibres.

\section*{Clinical Features}
The onset is usually abrupt, and symptoms may be at their worst within 48 hours.

\begin{itemize}
    \item \textbf{Facial weakness:} drooping of one side of the face, including the forehead, eyebrow, and corner of the mouth; the affected side may appear smooth and expressionless
    \item \textbf{Eye involvement:} difficulty closing the eye on the affected side, with dryness, irritation, and excessive watering
    \item \textbf{Oral symptoms:} drooling, difficulty eating and drinking, and slurred or unclear speech
    \item \textbf{Pain:} pain behind or in front of the ear on the affected side, often present at the onset
    \item \textbf{Taste disturbance:} a reduced or altered sense of taste on the front two-thirds of the tongue
    \item \textbf{Hearing symptoms:} sensitivity to loud sounds (hyperacusis) in the affected ear
    \item \textbf{Other sensory symptoms:} a dull or heavy feeling on the affected side of the face, although true numbness is unusual
\end{itemize}

A characteristic feature of Bell's palsy is that it affects the whole half of the face, including the forehead. Weakness that spares the forehead suggests a central (upper motor neurone) cause such as stroke.

\section*{Differential Diagnosis}
Any sudden facial weakness must be assessed urgently, because the most important condition to exclude is a stroke or transient ischaemic attack.

\begin{itemize}
    \item \textbf{Stroke or transient ischaemic attack:} a medical emergency; typically causes weakness of the lower half of the face that spares the forehead, and is often accompanied by limb weakness, speech problems, or visual disturbance
    \item \textbf{Ramsay Hunt syndrome:} shingles (varicella zoster) affecting the facial nerve, causing severe ear pain, a vesicular rash on or around the ear, and hearing loss
    \item \textbf{Lyme disease:} a bacterial infection spread by tick bites that can cause facial palsy, sometimes on both sides
    \item \textbf{Acoustic neuroma} (vestibular schwannoma) and other tumours of the facial nerve, parotid gland, or base of the skull
    \item \textbf{Sarcoidosis} and other inflammatory conditions affecting the facial nerve
    \item \textbf{Multiple sclerosis} and other demyelinating neurological conditions
    \item \textbf{Infections} such as otitis media, mastoiditis, or HIV
    \item \textbf{Trauma,} including fractures of the skull or temporal bone and surgical injury
\end{itemize}

\section*{When to Seek Medical Care}
Any sudden facial weakness should be assessed promptly so that stroke and other serious causes can be excluded. See a doctor urgently if facial weakness occurs together with:

\begin{itemize}
    \item Headache, fever, or a rash on or inside the ear
    \item Hearing loss, ringing in the ear, or dizziness
    \item Weakness or numbness spreading to an arm or leg
    \item Slurred speech, confusion, or difficulty swallowing
    \item Weakness affecting both sides of the face
    \item Inability to close the eye, which places the cornea at risk
\end{itemize}

\textbf{Contact your local emergency services for:}
\begin{itemize}
    \item Signs of stroke, remembered with the FAST mnemonic (Face, Arms, Speech, Time)
    \item Sudden loss of consciousness or a seizure
    \item Sudden severe headache or neck stiffness
    \item Difficulty breathing or swallowing
    \item Sudden weakness or numbness of the whole of one side of the body
\end{itemize}

\section*{Diagnosis and Investigations}
Bell's palsy is a clinical diagnosis, based on the pattern of weakness and a careful examination. There is no single confirmatory test.

\begin{itemize}
    \item \textbf{History:} timing and speed of onset, ear pain, preceding illness, rash, exposure to ticks, diabetes, and any previous episodes of facial weakness
    \item \textbf{Neurological examination:} assessment of all cranial nerves, including eye movement, hearing, and sensation, to confirm an isolated lower-motor-neurone facial nerve palsy
    \item \textbf{Ear examination:} inspection of the ear canal and tympanic membrane for vesicles, suggesting Ramsay Hunt syndrome, or signs of middle-ear infection
    \item \textbf{Blood tests:} used selectively for diabetes, inflammatory markers, Lyme disease serology, or HIV if indicated
    \item \textbf{Imaging:} MRI of the brain and internal auditory canals is not routine but is considered when the onset is slow, symptoms are atypical, there is progression beyond three weeks, or a tumour is suspected
    \item \textbf{Electromyography (EMG)} and nerve conduction studies: occasionally used later to assess the severity of nerve damage and to help predict recovery
\end{itemize}

\section*{Management}
\begin{itemize}
    \item \textbf{Corticosteroids:} prednisolone started within 72 hours of onset improves the chance of complete recovery and speeds the return of facial function; this is the mainstay of treatment
    \item \textbf{Antiviral medication:} such as aciclovir, is sometimes added when a viral cause such as shingles is suspected, although evidence that it adds benefit to steroids alone is limited
    \item \textbf{Eye protection:} lubricating eye drops during the day, ointment at night, and taping or patching the eyelid closed while sleeping to protect the cornea from drying and ulceration
    \item \textbf{Facial exercises and physiotherapy:} gentle massage and graded facial exercises may help maintain muscle tone and restore symmetry as function returns
    \item \textbf{Pain relief:} simple analgesics for ear or facial discomfort
    \item \textbf{Psychological support:} the sudden change in appearance can be distressing, and reassurance about the good prognosis is important
    \item \textbf{Follow-up:} review to monitor recovery and eye safety, with referral to neurology or ear, nose, and throat specialists for atypical, severe, or slow-to-recover cases; referral to an oculoplastics or ophthalmology service if the eye cannot be closed
\end{itemize}

\section*{Complications}
\begin{itemize}
    \item \textbf{Incomplete recovery:} residual facial weakness or asymmetry in a minority of people
    \item \textbf{Synkinesis:} abnormal, involuntary co-contraction of facial muscles, so that moving one part of the face produces movement elsewhere (for example the corner of the mouth moving when the eye is closed)
    \item \textbf{Eye complications:} dryness, corneal ulceration, and infection resulting from an eye that cannot be closed
    \item \textbf{Persistent facial spasm or contracture:} a tight, drawn appearance of the affected muscles
    \item \textbf{Crocodile tears:} excessive tearing of the affected eye while eating
    \item \textbf{Long-term sensory changes:} altered taste or hearing on the affected side
    \item \textbf{Psychosocial effects:} reduced self-confidence and distress related to facial appearance
\end{itemize}

\section*{Prognosis and Outcomes}
The outlook for Bell's palsy is generally excellent. Around 8 in 10 people recover substantially, with the most marked improvement occurring in the first three weeks and further recovery continuing over three to six months.

\begin{itemize}
    \item \textbf{Early treatment:} corticosteroids started within 72 hours of onset significantly improve the chance of complete recovery
    \item \textbf{Recovery rates:} without treatment, roughly 7 in 10 people recover completely; with early steroid therapy this rises to over 8 in 10
    \item \textbf{Persistent weakness:} around 15 per cent of people are left with some degree of permanent weakness or cosmetic abnormality
    \item \textbf{Predictors of worse outcome:} severe or complete paralysis at onset, older age, diabetes, and associated pain
    \item \textbf{Recurrence:} a small proportion of people experience a further episode, sometimes on the opposite side
\end{itemize}

\section*{Prevention and Self-Care}
Because the precise cause is not fully understood, Bell's palsy cannot be reliably prevented. During recovery, however, self-care is important.

\begin{itemize}
    \item Protect the affected eye from drying with preservative-free lubricating drops, and use ointment and eyelid taping at night
    \item Wear sunglasses or a protective eye shield outdoors to reduce irritation and dust exposure
    \item Eat and drink slowly, and use a straw if drinking is difficult
    \item Perform gentle facial massage and practise slow, deliberate facial movements as advised by a physiotherapist
    \item Attend medical review early so that corticosteroid treatment can start within the critical 72-hour window
    \item Maintain good general health, including control of diabetes and blood pressure, which are associated with a higher risk of the condition
\end{itemize}

\section*{Sources and Further Reading}
The following organisations provide reliable information on Bell's palsy for patients, families, and health professionals.

\begin{itemize}
    \item NHS. \textit{Information on Bell's palsy, including symptoms, treatment, and recovery.} https://www.nhs.uk/conditions/
    \item Mayo Clinic. \textit{Overview of Bell's palsy, its causes, and management.} https://www.mayoclinic.org/
    \item MSD Manual. \textit{Professional information on facial nerve palsy.} https://www.msdmanuals.com/
    \item MedlinePlus. \textit{Health information on Bell's palsy for patients and families.} https://medlineplus.gov/
    \item NICE. \textit{Clinical guidance on the diagnosis and management of Bell's palsy.} https://www.nice.org.uk/
\end{itemize}
